% Bagian: sets-functions-relations
% Bab: functions
% Subbagian: basics

\documentclass[../../../include/open-logic-section]{subfiles}

\begin{document}

\olfileid[id]{sfr}{fun}{bas}
\olsection{Dasar-Dasar}

\begin{explain}
Suatu \emph{fungsi} adalah pemetaan yang memetakan setiap !!{element} dari
suatu himpunan tertentu ke sebuah !!{element} tertentu dalam suatu himpunan
(lain) yang ditentukan. Sebagai contoh, operasi penambahan~$1$ menentukan
sebuah fungsi: setiap bilangan~$n$ dipetakan ke satu-satunya bilangan~$n+1$.

Secara lebih umum, fungsi dapat menerima pasangan, tripel, dan sebagainya
sebagai masukan serta menghasilkan suatu keluaran. Banyak fungsi sudah kita
kenal dari aritmetika dasar. Sebagai contoh, penjumlahan dan perkalian adalah
fungsi. Keduanya menerima dua bilangan dan menghasilkan bilangan ketiga.

Dalam pengertian matematis yang abstrak ini, fungsi adalah sebuah \emph{kotak
hitam}: yang penting hanyalah keluaran mana yang dipasangkan dengan masukan
mana, bukan cara menghitung keluaran tersebut.
\end{explain}

\begin{defn}[Fungsi]
Suatu \emph{fungsi} $f \colon A \to B$ adalah pemetaan setiap !!{element}
dari~$A$ ke sebuah !!{element} dari~$B$.

Kita menyebut $A$ sebagai \emph{domain} dari~$f$ dan $B$ sebagai
\emph{kodomain} dari~$f$. !!{element}s dari~$A$ disebut masukan atau
\emph{argumen} dari~$f$, dan !!{element} dari~$B$ yang dipasangkan dengan
argumen~$x$ oleh~$f$ disebut \emph{nilai dari~$f$} untuk argumen~$x$, yang
ditulis~$f(x)$.

\emph{Daerah hasil} $\ran{f}$ dari~$f$ adalah himpunan bagian kodomain yang
terdiri atas nilai-nilai~$f$ untuk suatu argumen; $\ran{f} =
\Setabs{f(x)}{x \in A}$.
\end{defn}

Diagram dalam \olref{fig:function} dapat membantu kita memahami fungsi. Elips
di sebelah kiri menyatakan \emph{domain} fungsi; elips di sebelah kanan
menyatakan \emph{kodomain} fungsi; dan sebuah panah menunjuk dari suatu
\emph{argumen} dalam domain ke \emph{nilai} yang bersesuaian dalam kodomain.

\begin{figure}
  \olasset{assets/diagrams/function.tikz}
  \caption{Suatu fungsi memetakan setiap !!{element} dari satu himpunan ke
    !!a{element} dari himpunan lain. Sebuah panah menunjuk dari suatu argumen
    dalam domain ke nilai yang bersesuaian dalam kodomain.}
  \ollabel{fig:function}
\end{figure}

\begin{ex}
Perkalian memetakan pasangan bilangan asli sebagai masukan ke bilangan asli
sebagai keluaran. Jadi, pemetaan itu berjalan dari $\Nat \times \Nat$ (domain)
ke $\Nat$ (kodomain). Ternyata daerah hasilnya juga
$\Nat$, karena setiap $n \in \Nat$ sama dengan $n \times 1$.
\end{ex}

\begin{ex}
Perkalian merupakan fungsi: setiap masukan, yakni pasangan bilangan asli,
dipasangkan dengan satu keluaran: $\times \colon \Nat^2 \to
\Nat$. Sebaliknya, pemetaan suatu bilangan asli~$n$ ke bilangan real~$x$ yang
memenuhi $x^2 = n$ bukanlah fungsi, karena setiap bilangan bulat positif~$n$
mempunyai dua akar kuadrat: $\sqrt{n}$ dan~$-\sqrt{n}$. Kita dapat menjadikannya
fungsi dengan hanya menghasilkan akar kuadrat utama, yaitu akar kuadrat nonnegatif:
$\sqrt{\phantom{X}} \colon \Nat \to \Real$.
\end{ex}

\begin{ex}
Relasi yang memasangkan setiap mahasiswa dalam suatu kelas dengan nilai akhir
merupakan fungsi; satu mahasiswa tidak dapat memperoleh dua nilai akhir yang
berbeda di kelas yang sama. Relasi yang memasangkan setiap mahasiswa dengan
orang tuanya bukanlah fungsi: seorang mahasiswa dapat
mempunyai nol, dua, atau lebih orang tua.
\end{ex}

\begin{explain}
Kita dapat mendefinisikan fungsi dengan menentukan secara saksama nilai fungsi
untuk setiap argumen yang mungkin. Beberapa caranya ialah dengan memberikan
rumus, menjelaskan metode untuk menghitung nilainya, atau mencantumkan nilai
untuk setiap argumen. Bagaimanapun suatu fungsi didefinisikan, kita harus
memastikan bahwa untuk setiap argumen ditentukan satu dan hanya satu nilai.
\end{explain}


\begin{ex}
Misalkan $f \colon \Nat \to \Nat$ didefinisikan oleh $f(x) = x+1$. Definisi
ini menentukan $f$ sebagai fungsi yang menerima bilangan asli dan menghasilkan
bilangan asli. Definisi ini menyatakan bahwa, jika diberi bilangan asli~$x$,
$f$ menghasilkan penerusnya~$x+1$.
Dalam hal ini, kodomain $\Nat$ bukan daerah hasil dari~$f$, karena bilangan
asli~$0$ bukan penerus bilangan asli mana pun. Daerah hasil dari~$f$ adalah
himpunan semua bilangan bulat positif, $\Int^{+}$.
\end{ex}

\begin{ex}\ollabel{examplefunext}
Misalkan $g \colon \Nat \to \Nat$ didefinisikan oleh $g(x) = x+2-1$.
Definisi ini menyatakan bahwa $g$ adalah fungsi yang menerima bilangan asli dan
menghasilkan bilangan asli. Jika diberi bilangan asli~$n$, $g$ menghasilkan
pendahulu dari penerus dari penerus bilangan masukan (yang dalam rumus
dilambangkan dengan~$x$), yaitu~$x+1$.
\end{ex}

\begin{explain}
Baru saja kita meninjau dua fungsi, $f$ dan $g$, dengan \emph{definisi} yang
berbeda. Namun, keduanya adalah \emph{fungsi yang sama}. Lagi pula, untuk
setiap bilangan asli~$n$, kita mempunyai $f(n) = n+1 = n+2-1 = g(n)$. Dengan
kata lain, definisi kita untuk $f$ dan~$g$ menentukan pemetaan yang sama melalui
persamaan yang berbeda. Jadi, secara tersirat kita mengandalkan prinsip
ekstensionalitas untuk fungsi,
\[
  \text{jika }\forall x\, f(x) = g(x)\text{, maka }f = g
\]
dengan syarat $f$ dan~$g$ mempunyai domain dan kodomain yang sama.
\end{explain}

\begin{ex}
Kita juga dapat mendefinisikan fungsi secara kasus. Sebagai contoh, kita dapat
mendefinisikan $h \colon \Nat \to \Nat$ dengan
\[
h(x) =
\begin{cases}
  \frac{x}{2} & \text{jika $x$ genap} \\
  \frac{x+1}{2} & \text{jika $x$ ganjil.}
\end{cases}
\]
Karena setiap bilangan asli itu genap atau ganjil, keluaran fungsi ini selalu
merupakan bilangan asli. Ingatlah bahwa jika Anda mendefinisikan fungsi secara
kasus, setiap masukan yang mungkin harus termasuk dalam tepat satu kasus.
Dalam keadaan tertentu, hal ini memerlukan pembuktian bahwa kasus-kasus
tersebut mencakup semua kemungkinan dan saling lepas.
\end{ex}

\end{document}
